\section{问题重述与分析}

\subsection{问题的统一结构}

玩家在图 (G=(V,E)) 上从起点走向终点。每天只能移动到相邻区域或停留；沙暴日不能移动；在矿山可用更高消耗换取收入，在村庄可用更高价格补给。水、食物、现金和负重共同限制可行性，到达终点后剩余物资按基准价一半回收。六关共享同一物理系统，但给定信息和玩家数量不同：

\begin{itemize}
  \item 问题一：未来天气完全已知，是确定性、有限时域的网络优化；
  \item 问题二：行动前只知当天天气，是满足非前视约束的反馈控制；
  \item 问题三：收益和消耗还依赖他人同日行动，是带拥挤外部性的动态博弈。
\end{itemize}

\begin{figure}[H]
  \centering
  \includegraphics[width=0.98\textwidth]{model_workflow.pdf}
  \caption{由信息集决定的统一模型与验证链}
  \label{fig:workflow}
\end{figure}

\subsection{四个建模难点}

\textbf{难点一是库存和位置强耦合。} 同一条几何最短路未必经济最优：晴朗、高温和沙暴的资源成本不同，绕行矿山可能增加现金，也可能因额外移动和负重失去可行性。

\textbf{难点二是信息时序。} 第三、四、六关只能依据日期、当前状态、当天气象和已观测历史行动，不得使用未来天气。若先生成整条天气再让策略读取，就是信息泄漏。

\textbf{难点三是“较少沙暴”并非概率。} 题面没有唯一概率分布，因而不存在脱离假设的唯一期望最优策略。本文把概率视为建模参数，分别报告低沙暴、基准和不利分布，并用风险预算筛选策略。

\textbf{难点四是多人规则带来负外部性。} 三人复制同一单人策略会在同一有向边形成 (2k) 倍消耗，在矿山又分摊收益。问题三必须同时分析路线错峰、矿山竞争、角色公平和单边偏离。

\subsection{地图规范化}

附件地图被重建为无向邻接表，并自动检查节点完备、邻接对称和全图连通。第三/五关、第四/六关分别共享同一图，因此六关实际只涉及四种图结构，见\cref{fig:maps}。图中第一、二关的绿线是求解器输出的最优轨迹，不是人工描线。

\begin{figure}[H]
  \centering
  \includegraphics[width=0.98\textwidth]{maps_overview.pdf}
  \caption{六关地图的规范图结构与确定天气最优轨迹}
  \label{fig:maps}
\end{figure}
